\section{统计分析、校准与灵敏度}\label{app:statistics}\label{app:diagnostics}
本节给出正文使用的留出重复诊断。实证估计使用共同的 $\RcvPairedN$ 道完整 MATH 题；合成校准采用单独指定的样本量。

\subsection{留出重复统计量与判定规则}
三折中的每一折，用两次发现重复为每题选择一个成员，并按平均发现准确率选择一个最佳固定成员；第三次重复将逐题映射对照该固定成员评分。逐题增益在折上平均得到 $G$，两臂的配对差为 $D=G_{\rm real}-G_{\rm clone}$（\S\ref{sec:calibration-test}）。重复平均 oracle 插件量 $\widehat H$ 仅作描述，不决定判定结果。

发现分数落在 $\{0,\tfrac12,1\}$ 上，因此并列很常见。选择使用 iid $U(0,10^{-9})$ 抖动对称地打破成员并列；实证报告对十个冻结的并列实现取平均，每个校准数据集使用一个并列实现。实证随机化参考在每臂的每题、每次重复内独立置换成员标签。每次抽取都重算两个选择、两个增益及 $D$，共 $10{,}000$ 次，并使用 $+1$ Monte Carlo 校正。配对题目 bootstrap 使用固定随机种子的 $2{,}000$ 次抽取，两臂共享重采样题目身份，以获得 95\% 区间。

正向支持要求同时满足：在完整案例筛选之前，任一臂及两臂交集的不完整题目比例不超过 $10\%$；配对题目至少 $30$ 道；随机化 $p<0.05$；$D$ 的 bootstrap 下界为正；$D>1.0$\,pp；且 $G_{\rm real}$ 的下界为正。当前比较为 \textsc{Insufficient Evidence}，不确立等效性。执行条件与程序身份检查要求匹配且独立获取的重复。仅凭结果矩阵相同不能认定重复采集，仍需核实不同的执行来源。

\paragraph{分析修订。}
采集协议先于采集制定；当前 v3 诊断与按完整性筛选的分析在采集后确定。原停止规则适用于初始采集阶段。补采集与共同集重分析作为独立修订记录，当前估计绑定 9 月 26 日快照。附录~\ref{app:wp1r} 给出纳入规则、排除题特征、评分冲突与冲突排除敏感性。更早的诊断版本已被替代，不用于当前估计。

\subsection{总体目标与诊断边界}
对于成员成功概率 $q_h(x)$，稳定互补性为
\[
H_{\rm stable}=\mathbb E_x\max_h q_h(x)-\max_h\mathbb E_x q_h(x).
\]
当某个固定成员在几乎每题上都是期望最佳时，该量为零。选择器相对于\emph{总体最佳固定}成员的总体增益 $\theta^*$ 满足 $\theta^*\le H_{\rm stable}$，因为其逐题成功概率不超过 $\max_h q_h(x)$。

实证 $G$ 的比较对象由有限发现数据选出，未必是总体最佳固定成员。克隆相减提供观测采样参照，并不解析消除全部有限样本选择偏差。因此，所报告的 $D$ 区间不是 $H_{\rm stable}$ 的区间；接近零的 $D$ 不能把稳定互补性界定在零附近。

在成员标签可交换的等能力边界处，随机化参考是精确的。对于更广的全局支配零假设，错误控制仅在校准网格上经验检查，没有一般有限样本定理保证。

\paragraph{有限重复的 oracle 估计。}
对重复均值 $\hat q_h(x)$，Jensen 不等式逐题给出 $\mathbb E\max_h\hat q_h(x)\ge\max_h\mathbb E\hat q_h(x)$。余量还减去第二个最大值，因此该不等式一般并不意味着二者之差向上偏。在 $q=0.9$、$M=9$、$T=400$、$R=3$ 的已知真值等能力模拟中，$300$ 个数据集的平均插件量为 $8.69$\,pp，而 $H_{\rm stable}=0$。实证同代码臂同样给出 $\RcvClonePlugin$\,pp。这些例子表明，正插件余量不足以证明专长。

\subsection{已知真值校准与统计功效}
二元张量从指定的成员—任务概率矩阵抽取。零假设设置包括等能力、共享题目难度、全局支配、$0.5/0.9$ 支配配置、异质支配，以及匹配的黏性题内潜在状态。备择包括潜在任务类型间的交叉与难度相关的部分交叉。网格覆盖 $M\in\{3,9,20\}$、$T\in\{100,400,1000\}$、$R\in\{3,5,10,20\}$、完整数据、5\% MCAR 与 3\% 难度相关缺失。

表~\ref{tab:calibration} 保留 $M=9,T=400$ 的设置，每格使用 $300$ 个 Monte Carlo 数据集（扫描中为 $150$ 个），每个数据集执行 $300$ 次随机化，比较插件量与当前克隆校准规则的支持率。在所测 $100$ 格网格上，该规则最坏的观测零假设错误支持率为 $2\%$。真实 $H_{\rm stable}=17.5$\,pp 的双专家备择在 $R=3$ 时功效为 $0\%$，在 $R=5$ 时约为 $0.7\%$，在 $R=10,20$ 时为 $100\%$。更小的循环交叉即使在 $R=20$ 时仍大多无法检出。这些设置特定的比率不构成一般检验大小或功效保证；当前研究的三次重复仍未判清稳定互补性。

\begin{table}[H]
\centering\small
\renewcommand{\arraystretch}{0.95}
\setlength{\tabcolsep}{4pt}
\footnotesize
\begin{tabular}{@{}lcccccccc@{}}
\toprule
 & \multicolumn{2}{c}{$R{=}3$} & \multicolumn{2}{c}{$R{=}10$} & \multicolumn{2}{c}{$R{=}20$} & abst. \\
\cmidrule(lr){2-3}\cmidrule(lr){4-5}\cmidrule(lr){6-7}\cmidrule(lr){8-8}
regime & plug & v3 & plug & v3 & plug & v3 & $R{=}3$ \\
\midrule
equal $q{=}0.9$$^\dagger$ & 100 & 0 & 100 & 0 & 100 & 1 & 100 \\
equal + difficulty$^\dagger$ & 100 & 0 & 100 & 1 & 100 & 1 & 100 \\
$0.5/0.9$ dominance$^\dagger$ & 100 & 0 & 100 & 0 & 13 & 0 & 100 \\
dominance$^\dagger$ & 100 & 0 & 100 & 0 & 100 & 0 & 100 \\
cross $H{\approx}0.9$\,pp & 100 & 1 & 100 & 0 & 100 & 0 & 99 \\
cross $H{\approx}1.8$\,pp & 100 & 1 & 100 & 1 & 100 & 0 & 99 \\
cross $H{\approx}4.2$\,pp & 100 & 0 & 100 & 0 & 100 & 0 & 100 \\
two strong specialists, $H{=}17.5$\,pp & 100 & 0 & 100 & 100 & 100 & 100 & 100 \\
\bottomrule
\end{tabular}

{\scriptsize $^\dagger$Null ($H{=}0$): false-support \%, nominal 5\%; else power \%. $M{=}9,T{=}400$, $B{=}300$ at $R{=}3$ ($150$ sweeps), 300 randomizations. Wilson intervals and the full $M,T,R$, missingness grid in Appendix~\ref{app:diagnostics}.}

\caption{\textbf{三次重复甚至会漏掉强正例。}在 $M=9,T=400$ 下，零假设行报告错误支持率（\%），备择行报告功效（\%）；最后一列为 $R=3$ 的弃权率。克隆校准规则（v3）很少支持零假设，但对所测小、中等交叉的灵敏度很低。}
\label{tab:calibration}
\end{table}
